% Chapter 10: Integration — the unified platform

\chapter{Integration: The Unified Platform}

The six papers presented in this thesis share a common Python package
(\texttt{satellite-paraguay}) that enables efficient execution, cross-paper
insights, and a unified deployment platform.

\section{Shared Infrastructure}

All six papers use the same:

\begin{itemize}
    \item Paraguay admin boundary data (18 deptos, 268 distritos, 7,912 tiles)
    \item Sentinel-2 imagery ingestion pipeline
    \item MapBiomas + Hansen ground truth
    \item OpenAQ + Sentinel-5P atmospheric data
    \item Catastro + indigenous territory conflict detection
    \item MLflow experiment tracking
    \item FastAPI endpoints
    \item Streamlit dashboard
\end{itemize}

This shared infrastructure reduces duplication and enables efficient execution.
Each paper's specific additions are in \texttt{src/papers/PXXXX\_name/}.

\section{Cross-Paper Insights}

Combining the six papers yields insights that no individual paper could:

\begin{enumerate}
    \item \textbf{Deforestation (P0011) + Carbon Credits (P0100):} Deforested
    areas can be cross-referenced with Verra projects to verify claims.
    \item \textbf{Yield (P0025) + Air Quality (P0035):} Crop residue burning
    (a major P0025 issue) drives PM2.5 spikes in Asunción (P0035 finding).
    \item \textbf{Indigenous (P0012) + Deforestation (P0011):} Indigenous
    territories face higher deforestation pressure than non-indigenous areas
    \citep{fao2024paraguay}.
    \item \textbf{Wildlife (P0026) + FIRMS:} NASA FIRMS fire alerts can be
    used to find poaching camps (validated via P0026 satellite imagery).
\end{enumerate}

\section{Deployment Architecture}

The package is deployed in three ways:

\begin{enumerate}
    \item \textbf{Command-line:} \texttt{make run-paper-1} etc.
    \item \textbf{API:} \texttt{make api} → FastAPI on port 8000
    \item \textbf{Dashboard:} \texttt{make dashboard} → Streamlit on port 8501
\end{enumerate}

Docker support is included for all three.

\section{Stakeholder Integration}

The thesis engages with multiple stakeholders (see \texttt{docs/STAKEHOLDERS.md}):

\begin{itemize}
    \item Cristaldo Lab (FADA-UNA) — primary advisor
    \item INFONA — forestry data (P0011, P0100)
    \item INDI — indigenous territories (P0012)
    \item SENEPA — public health (P0031)
    \item UN-Habitat — indigenous atlas (P0012)
    \item WWF Paraguay, Guyra — wildlife (P0026)
\end{itemize}

\section{Open-Source Contribution}

The complete code, data catalog, models, and documentation are released
under open-source licenses. We expect this to support future Paraguayan
researchers and serve as a template for similar efforts in other Latin
American countries.

\section{Lessons Learned}

\begin{itemize}
    \item Foundation models can be fine-tuned for Paraguay with limited data
    \item The CARE Principles framework is essential for indigenous data work
    \item Multi-modal approaches (vision + language + time-series) are powerful
    \item Local data (paraguay-geodata) is critical for Paraguay-specific work
    \item Open data + open models enables world-class research at low cost
\end{itemize}
